\subsection{Polynomial rings over a factorial domain}

Let A be a factorial domain, K its field of fractions and f a non-zero element of K{[}X{]}; an element c of K* will be called a content of f if it is a g.c.d. of the coefficients of f.~Let v be a valuation on K which is essential for A and \(\bar{v}\) its canonical extension to K{[}X{]} (defined by \(\bar{v}\left(\sum_{i}a_{i}X^{i}\right)=\inf v(a_{i})\) ; cf.~Chapter VI, § 10, no. 1, Proposition 2); then \(\bar{v}(f)=v(c)\) .

LEMMA 1 (Gauss). Let \(f, f'\) be non-zero elements of \(\mathbf{K}[\mathbf{X}]\) and \(c, c'\) contents of \(f, f'\) . Then \(cc'\) is a content of \(ff'\) .

Let \(d\) be a content of \(ff'\) . For every valuation \(v\) on \(K\) which is essential for \(A\) , let \(\bar{v}\) denote its canonical extension to \(K[X]\) . Then

\[
v (d) = \bar {v} (f f ^ {\prime}) = \bar {v} (f) + \bar {v} (f ^ {\prime}) = v (c) + v (c ^ {\prime}) = v (c c ^ {\prime}).
\]

Hence \(cc'd^{-1}\) is an invertible element of A.

THEOREM 2. Let A be a factorial domain, K its field of fractions, \((p_{\iota})\) a representative system of extremal elements of A and \((\mathrm{P}_{\lambda})\) a representative system of irreducible polynomials of K{[}X{]}, each \(P_{\lambda}\) having 1 as a content. Then:

\begin{enumerate}
\def\labelenumi{(\roman{enumi})}
\item
  A{[}X{]} is a factorial domain;
\item
  the set of \(p_{\iota}\) and \(P_{\lambda}\) is a representative system of extremal elements of A{[}X{]}.
\end{enumerate}

Let \(f\) be a non-zero element of \(A[X]\) . In the ring \(K[X]f\) can be decomposed uniquely in the form:

\[
f = a \prod_ {\lambda} \mathrm{P} _ {\lambda} ^ {n (\lambda)} \quad (a \in \mathrm{K} ^ {*}, n (\lambda) \geqslant 0).
\]

Lemma 1 proves that \(a\) is a content of \(f\) . Hence \(a \in \mathbf{A}\) . As \(\mathbf{A}\) is factorial, \(a\) can be decomposed uniquely in the form:

\[
a = u \prod_ {\iota} p _ {\iota} ^ {m (\iota)} \quad (u \text {   invertible   in   A }, m (\iota) \geqslant 0).
\]

Whence the existence and uniqueness of the decomposition:

\[
f = u \prod_ {\iota} p _ {\iota} ^ {m (\iota)} \prod_ {\lambda} \mathrm{P} _ {\lambda} ^ {n (\lambda)}.
\]

Note that this theorem proves that every element of A admits the same decomposition into extremal elements in A and A{[}X{]}. The g.c.d. of a family of elements of A is therefore the same in A and in A{[}X{]}.

We may also use Proposition 18 of § 1, no. 10 to show that A{[}X{]} is a factorial domain if and only if A is a factorial domain.

COROLLARY. If A is a factorial domain, the domain A{[}X \(_{1}\) , . . . , X \(_{n}\) {]} is factorial.

Argue by induction on n.

This corollary may be extended to the case of an infinite family of indeterminates (cf.~Exercise 2).

